# 化学基础知识

## 气体

气体基本特性：

1. 自动扩散性
2. 可压缩性

### 气体的状态方程

#### 理想气体的状态方程

假设分子微粒之间的作用力、分子本身的体积忽略不计的气体为**理想气体**。

（**低压高温**状态）理想气体状态方程 $pV=nRT$

#### 实际气体的状态方程

Van der Waals 气体状态方程 $\left(p+a\left({\dfrac{n}{V}}\right)^2\right)(V-bn)=nRT$，其中 $p_内=\left(\dfrac{n_外}{V}\right)\left(\dfrac{n_内}{V}\right)=\left(\dfrac{n}{V}\right)^2$，$a$ 是与分子间引力有关的常数，$b$ 是与分子本身体积有关的常数，统称为 van der Waals 常数。通常，$a$ 和 $b$ 数值越大，表明实际气体距理想气体愈远，从而产生的偏差就愈大。

### 混合气体的分压定律

Dalton 分压定律 $\displaystyle p_总=\sum_ip_i$
$p_i=p_总x_i$

Amagat 分体积定律 $\displaystyle V_总=\sum_iV_i$

### 气体扩散定律

$v \propto \sqrt{\dfrac{1}{\rho}}$
$v \propto \sqrt{\dfrac{1}{M}}$

### 气体分子的速率分布和能量分布

#### 气体分子的速率分布

1859 年，英国物理学家 Maxwell 用概率论的方法推导出了计算 气体分子运动速率分布的公式，讨论了分子运动速率的分布。后来，奥地利人 L. Boltzmann 用统计力学的方法也得到了相同的结果，从而进一步验证了麦克斯韦的速率分布公式。

#### 气体分子的能量分布

在无机化学中，常用下面的近似公式来计算和讨论能量的分布。

$$
f_{E_0}=\dfrac{N_i}{N}=\exp{\dfrac{-E_0}{RT}}
$$

## 液体和溶液

**溶液**是指一种或一种以上的物质以分子或离子的形式溶解在另一种物质中形成的均一 、稳定的混合物。其中物质的量多且能溶解其他物质的组分称为**溶剂**，含量少的组分称为**溶质**。

### 溶液浓度的表示方法

对于溶剂 $\rm A$ 和溶质 $\rm B$，有

1. 质量摩尔浓度 $b({\rm B})=\dfrac{n({\rm B})}{m({\rm A})}$
2. 物质的量浓度 $c({\rm B})=\dfrac{n({\rm B})}{V}$
3. 质量分数 $w({\rm B})=\dfrac{m({\rm B})}{m}$
4. 摩尔分数 $x({\rm B})=\dfrac{n({\rm B})}{n}$

$x\approx \dfrac{1000}{18}b~{\rm mol^{-1}}$

### 溶液的饱和蒸气压

#### 纯溶剂的饱和蒸气压

当凝聚速率和蒸发速率相等时，饱和蒸气所具有的压强叫做该温度下液体的**饱和蒸气压**，简称**蒸气压**，用 $p^*$ 表示。对同一液体来说，其饱和蒸气压随温度的升高而增大。

#### 溶液的饱和蒸气压下降

溶液的饱和蒸气压 $p$ 小于纯溶剂的蒸气压 $p^*$。
溶液的浓度越大，溶质的粒子数目越多，溶液的饱和蒸气压比纯溶剂的饱和蒸气压下降得就越多。

Raoult 定律 $p=p^*x({\rm A})$

$\begin{align}\Delta p &= p^*-p\\&=p^*[1-x({\rm A})]\\&=p^*x({\rm B})\end{align}$

稀溶液的饱和蒸气压下降值与稀溶液的质量摩尔浓度成正比。

### 稀溶液的依数性

[稀溶液的饱和蒸气压下降](#溶液的饱和蒸气压下降)是稀溶液的一种依数性质。

#### 溶液的沸点升高

液体汽化的方式有两种，即**蒸发**和**沸腾**。 当液体的饱和蒸气压与外界大气压强相等时，液体的汽化将在其表面和内部同时发生，称为**沸腾**，这时的温度称为该液体的**沸点**。 而在低于此温度下的汽化，仅在液体表面上进行，称为**蒸发**。

难挥发性非电解质稀溶液沸点升高值 $\Delta T_\mathrm{b}=T_\mathrm{b}-T_\mathrm{b}^*$，与其饱和蒸气压下降的数值成正比，即 $\Delta T_\mathrm{b} \propto \Delta p$，记 $k_\mathrm{b}$ 为稀溶液**沸点升高常数**，得 $\Delta T_\mathrm{b}=k_\mathrm{b}b$（**沸点升高公式**）

#### 溶液的凝固点降低

液体凝固成固体（严格说是晶体）时的温度称为该液体的**凝固点**。由于溶液的蒸气压下降导致溶液的凝固点低于纯溶剂的凝固点。

难挥发性非电解质稀溶液凝固点降低值 $\Delta T_\mathrm{f}=T_\mathrm{f}^*-T_\mathrm{f}$，与其饱和蒸气压下降的数值成正比，记 $k_\mathrm{f}$ 为**凝固点降低常数**，得 $\Delta T_\mathrm{f}=k_\mathrm{f}b$（凝固点降低公式）

#### 渗透压

溶剂分子透过半透膜进入溶液的现象，称为**渗透现象**。渗透平衡时两侧液面高度差造成的静压称为溶液的**渗透压**。

Van’t Hoff 指出 $\mathit{\Pi}V = nRT$

## 固体和晶体

### 晶体和非晶体

#### 非晶体的特性

#### 晶体的特性

### 对称性

#### 旋转和旋转轴

#### 反映和镜面

#### 反演和对称中心

### 晶体和点阵

### 晶系和点阵型式

#### 七个晶系

#### 十四种空间点阵型式

### 晶胞
